% Part: proof-theory
% Chapter: natural-deduction
% Section: rules-N1

\documentclass[../../../include/open-logic-section]{subfiles}

\begin{document}

\begin{table}
\begin{center}
\begin{tabular}[t]{@{}cc@{}}
\bottomAlignProof
\AxiomC{$\Discharge{!A}{x}$}
\shortDeduce
\DeduceC{$!B$}
\DischargeRule{\Intro{\lif}}{x}
\UnaryInfC{$!A \lif !B$}
\DisplayProof
&
\bottomAlignProof
\AxiomC{$!A \lif !B$}
\AxiomC{$!A$}
\RightLabel{\Elim{\lif}}
\BinaryInfC{$!B$}
\DisplayProof
\\[4ex]
\bottomAlignProof
\AxiomC{$!A$}
\AxiomC{$!B$}
\RightLabel{\Intro{\land}}
\BinaryInfC{$!A \land !B$}
\DisplayProof
&
\begin{tabular}[b]{@{}l@{}}
\AxiomC{$!A \land !B$}
\RightLabel{\Elim{\land}[1]}
\UnaryInfC{$!A$}
\DisplayProof
\\[3ex]
\bottomAlignProof
\AxiomC{$!A \land !B$}
\RightLabel{\Elim{\land}[2]}
\UnaryInfC{$!B$}
\DisplayProof
\end{tabular}
\\[4ex]
\begin{tabular}{@{}l@{}}
\AxiomC{$!A$}
\RightLabel{\Intro{\lor}[1]}
\UnaryInfC{$!A \lor !B$}
\DisplayProof
\\[3ex]
\AxiomC{$!B$}
\RightLabel{\Intro{\lor}[2]}
\UnaryInfC{$!A \lor !B$}
\DisplayProof
\end{tabular}
&
\AxiomC{$!A \lor !B$}
\AxiomC{$\Discharge{!A}{x}$}
\shortDeduce
\DeduceC{$!C$}
\AxiomC{$\Discharge{!B}{y}$}
\shortDeduce
\DeduceC{$!C$}
\DischargeRule{\Elim{\lor}}{x\,y}
\TrinaryInfC{$!C$}
\DisplayProof
\bottomAlignProof
\\[6ex]
\bottomAlignProof
\AxiomC{$\Discharge{!A}{x}$}
\shortDeduce
\DeduceC{$\lfalse$}
\DischargeRule{\Intro{\lnot}}{x}
\UnaryInfC{$\lnot !A$}
\DisplayProof
&
\bottomAlignProof
\AxiomC{$\lnot !A$}
\AxiomC{$!A$}
\RightLabel{\Elim{\lnot}}
\BinaryInfC{$\lfalse$}
\DisplayProof
\\[4ex]
\bottomAlignProof
\AxiomC{$!A(a)$}
\RightLabel{\Intro{\lforall}}
\UnaryInfC{$\lforall[x][!A(x)]$}
\DisplayProof
&
\bottomAlignProof
\AxiomC{$\lforall[x][!A(x)]$}
\RightLabel{\Elim{\lforall}}
\UnaryInfC{$!A(t)$}
\DisplayProof
\\[4ex]
\bottomAlignProof
\AxiomC{$!A(t)$}
\RightLabel{\Intro{\lexists}}
\UnaryInfC{$\lexists[x][!A(x)]$}
\DisplayProof
&
\bottomAlignProof
\AxiomC{$\lexists[x][!A(x)]$}
\AxiomC{$\Discharge{!A(a)}{x}$}
\shortDeduce
\DeduceC{$!C$}
\DischargeRule{\Elim{\lexists}}{x}
\BinaryInfC{$!C$}
\DisplayProof
\\[4ex]
\bottomAlignProof
\AxiomC{$\lfalse$}
\RightLabel{\FalseInt}
\UnaryInfC{$!A$}
\DisplayProof
&
\bottomAlignProof
\AxiomC{$\Discharge{\lnot!A}{x}$}
\shortDeduce
\DeduceC{$\lfalse$}
\DischargeRule{\FalseCl}{x}
\UnaryInfC{$!A$}
\DisplayProof
\end{tabular}
\end{center}
\caption{\Log{N1i}, \Log{N1c} నియమాలు. $\Intro{\lforall}$లో \tetoken{స్థిరాంకం}{!!{constant}}~$a$ ఫలితంలో గానీ తెరిచి ఉన్న ఉపపత్తిలో గానీ ఉండకూడదు. $\Elim{\lexists}$లో $a$ ఫలితంలో గానీ ఆ నియమం తరువాత తెరిచి మిగిలే ఏ ఉపపత్తిలో గానీ ఉండకూడదు. \Log{N1i}లో \FalseCl ఉండదు.}
\ollabel{tab:N1}
\end{table}
\sourcecorrection{OLTEPTNATN1-001}{మూల శీర్షికలో $\Elim{\lexists}$ ఐగెన్ స్థిరాంకం పూర్వపక్షాలలో ఉండకూడదని చెప్పింది; అయితే ఉపపూర్వపక్షం $!A(a)$లో అది తప్పనిసరిగా ఉంది. ఫలితంలో, తెరిచి మిగిలే ఉపపత్తుల్లో ఉండకూడదని సరిచేశాం.}

\end{document}
